\section{Application to Tao's theorem}\label{sec:consequences}

Take $D=\Opos$, $F=\Syr$, and $\alpha=1001/1000$.  Let
\[
 R_s=\Opos\cap[s,s^\alpha],\qquad
 h_s=\sum_{n\in R_s}\frac1n,\qquad
 \mu_s(\{n\})=\frac{1}{h_s n}\quad(n\in R_s),
\]
whenever $R_s\ne\varnothing$.  The measure is zero outside $R_s$.
Choose $s_*$ sufficiently large that all blocks used below are nonempty.

\begin{theorem}[Tao, Proposition 1.11]\label{thm:tao-input}
There are absolute constants $B,c>0$ such that, for all sufficiently
large $x$ and $y\in\{x^\alpha,x^{\alpha^2}\}$, an input $\mathbf N_y$
with law $\mu_y$ satisfies
\[
 \Prob(t_x(\mathbf N_y)=\infty)\leq Bx^{-c}.
\]
Let $\widehat p_x(n)$ agree with $p_x(n)$ on successful passage and
equal $1$ on failed passage, as in Tao's convention.  Then
\begin{equation}\label{eq:tao-tv}
 \normone{(\widehat p_x)_*\mu_{x^\alpha}
 -(\widehat p_x)_*\mu_{x^{\alpha^2}}}
 \leq B(\log x)^{-c}.
\end{equation}
\end{theorem}

The norm in \eqref{eq:tao-tv} is the unhalved $\ell^1$ norm used in
\cite{Tao7}.  Tao proves this proposition using the valuation law and
fine-scale mixing of Syracuse random variables, obtained through Fourier
estimates and a two-dimensional renewal argument.  The estimates in
Section~\ref{sec:repairs} supply the endpoint margin and uniform
coefficient bound used in his Section~5.

The change to $\dd$ costs only the known failure probability.  Indeed,
if $q$ merges $\dd$ with $1$ and fixes the other points, then for two
probability measures $\lambda,\lambda'$ on $\Opos\cup\{\dd\}$,
\begin{equation}\label{eq:unmerge}
 \normone{\lambda-\lambda'}
 \leq\normone{q_*\lambda-q_*\lambda'}
 +2|\lambda(\dd)-\lambda'(\dd)|.
\end{equation}
To check it, only the coordinates $1$ and $\dd$ need attention:
$|a|+|b|\leq|a+b|+2|b|$ for their respective differences.
By \eqref{eq:kernel},
$K_x\nu_{x^\alpha}=K_x\mu_{x^{\alpha^2}}$.
Thus Theorem~\ref{thm:tao-input} proves, after increasing a constant,
\begin{equation}\label{eq:input-completed}
 d(s)\leq Bs^{-c},\qquad e(s)\leq B'(\log s)^{-c}
 \quad(s\geq s_*).
\end{equation}
These measures retain their failure mass and the original denominators $h_s$.

\subsection{An exactly compatible family of first-entry laws}

Fix one base $b\geq s_*$ and put $t_j=b^{\alpha^j}$ for $j\geq0$.
For every real threshold $x\geq1$, write
\[
 E_x=(\Opos\cap[1,x])\cup\{\dd\},\qquad
 L_c=\frac{B'}{1-\alpha^{-c}},
\]
and define the probability measure
\begin{equation}\label{eq:lattice-law}
 \lambda_{x,j}^{(b)}=K_x\mu_{t_{j+1}}\quad\hbox{on }E_x.
\end{equation}
The input sequence is the same for every threshold. These laws describe
the first entry into $[1,x]$, including failure, rather than the position
after a prescribed number of iterations.

\begin{corollary}[Coherent first-entry limits]\label{cor:limit}
For each $x\geq1$, the measures $\lambda_{x,j}^{(b)}$ converge in total
variation to a probability measure $\lambda_x^{\infty,(b)}$ on $E_x$.
For every $j\geq0$ and $1\leq x\leq t_j$,
\begin{equation}\label{eq:lattice-rate}
 \normone{\lambda_{x,j}^{(b)}-\lambda_x^{\infty,(b)}}
 \leq L_c(\log t_j)^{-c}.
\end{equation}
The limiting family is exactly compatible at every pair of thresholds:
\begin{equation}\label{eq:lattice-compatible}
 K_x\lambda_y^{\infty,(b)}=\lambda_x^{\infty,(b)}
 \qquad(1\leq x\leq y).
\end{equation}
For $x\geq b$, its failure mass obeys
\begin{equation}\label{eq:all-threshold-failure}
 \lambda_x^{\infty,(b)}(\dd)
 \leq Bx^{-c/\alpha}+L_c\alpha^c(\log x)^{-c}.
\end{equation}
\end{corollary}

\begin{proof}
When $x\leq t_j$, nested passage gives
\[
 \lambda_{x,j}^{(b)}=K_x\nu_{t_j},\qquad
 \lambda_{x,j+1}^{(b)}=K_xK_{t_j}\nu_{t_{j+1}}.
\]
Hence contraction and \eqref{eq:input-completed} imply
\[
 \normone{\lambda_{x,j+1}^{(b)}-\lambda_{x,j}^{(b)}}
 \leq e(t_j)\leq B'(\alpha^j\log b)^{-c}.
\]
For each fixed $x$, these inequalities apply from some index onward.
The geometric tail makes the sequence Cauchy in the finite-dimensional
space $\ell^1(E_x)$. Its limit is nonnegative with total mass one.
Summing the tail from $j$ proves \eqref{eq:lattice-rate}, uniformly for
all $x\leq t_j$.

For $x\leq y$, the identity
$K_x\lambda_{y,j}^{(b)}=\lambda_{x,j}^{(b)}$ holds at every finite $j$.
Taking limits under the contraction $K_x$ proves
\eqref{eq:lattice-compatible}. In particular, the thresholds need not
belong to the sequence $(t_j)$.

For $x\geq b$, choose the unique $k\geq0$ with
$r=t_k\leq x<t_{k+1}=r^\alpha$. Then $r>x^{1/\alpha}$, and
$\lambda_{r,k}^{(b)}=\nu_r$. Equations \eqref{eq:input-completed} and
\eqref{eq:lattice-rate} give
\[
 \lambda_r^{\infty,(b)}(\dd)\leq Br^{-c}+L_c(\log r)^{-c}.
\]
Since $p_r(\dd)=\dd$, compatibility gives
$\lambda_x^{\infty,(b)}(\dd)\leq\lambda_r^{\infty,(b)}(\dd)$.
Substitution of $r>x^{1/\alpha}$ proves
\eqref{eq:all-threshold-failure}.
\end{proof}

The base $b$ records the logarithmic scale phase. Replacing it by
$b^{\alpha^k}$, $k\geq0$, deletes finitely many initial inputs and
therefore gives exactly the same limiting family. Equality for arbitrary
different bases is not supplied by this argument. For a threshold
$x=t_k$, the limit also equals the threshold-based limit obtained from
inputs $\mu_{x^{\alpha^{j+1}}}$, since that sequence is a tail of the
common input sequence.

\begin{corollary}[Joint first-entry distributions]\label{cor:joint-entry}
Fix $1\leq x_1\leq\cdots\leq x_r$, with $r\geq1$, and let
$\mathbf N_j$ have law $\mu_{t_{j+1}}$. The joint law of
\[
 (p_{x_1}(\mathbf N_j),\ldots,p_{x_r}(\mathbf N_j))
\]
converges to a probability measure $\Lambda_{x_1,\ldots,x_r}^{(b)}$.
It is supported on the finite set
\[
 \mathcal C=\{(z_1,\ldots,z_r)\in\textstyle\prod_i E_{x_i}:
                p_{x_i}(z_{i+1})=z_i\ (1\leq i<r)\}.
\]
For every $j$ with $t_j\geq x_r$, its unhalved total-variation error
is at most $L_c(\log t_j)^{-c}$, with no factor depending on $r$.
The marginal at $x_i$ is $\lambda_{x_i}^{\infty,(b)}$.
\end{corollary}

\begin{proof}
Define
\[
 J:E_{x_r}\longrightarrow\mathcal C,\qquad
 J(z)=(p_{x_1}(z),\ldots,p_{x_r}(z)).
\]
Nested passage proves that $J$ takes values in $\mathcal C$. Its last
coordinate is $z$, so it is injective. Conversely, repeated use of the
defining equations of $\mathcal C$ gives $z_i=p_{x_i}(z_r)$; thus the
inverse of $J$ is projection to the last coordinate. In particular $J$
has singleton fibres and loses no information. Its pushforward preserves
the $\ell^1$ norm of signed measures, because distinct source atoms have
distinct images.

The finite joint law is exactly $J_*\lambda_{x_r,j}^{(b)}$.
Set $\Lambda_{x_1,\ldots,x_r}^{(b)}=J_*\lambda_{x_r}^{\infty,(b)}$.
The isometry and \eqref{eq:lattice-rate} give the stated error bound.
Projection onto coordinate $i$ and
\eqref{eq:lattice-compatible} identify the marginal.
\end{proof}

Tao's Proposition~1.11 compares two successive starting scales. The
corollaries above turn its summable error into exact compatible
distributions, including simultaneous first-entry observations. The
compatibility maps are the first-passage maps $K_x$; no invariance under
one application of $\Syr$ is required.

\subsection{The quantitative fixed-threshold theorem}

\begin{theorem}[Tao's fixed-threshold theorem]
\label{thm:fixed}
There are absolute constants $C,c>0$ such that, for all real $K,X\geq2$,
\begin{align}
 \sum_{\substack{m\leq X,\ m\in\Opos\\\Syr_{\min}(m)>K}}\frac1m
 &\leq C\frac{\log X}{(\log K)^c},\label{eq:odd-fixed}\\
 \sum_{\substack{n\leq X\\\Col_{\min}(n)>K}}\frac1n
 &\leq C\frac{\log X}{(\log K)^c}.\label{eq:full-fixed}
\end{align}
\end{theorem}
\begin{proof}
Equations \eqref{eq:input-completed} and \eqref{eq:power-error}, followed
by Theorem~\ref{thm:transport}, give
\[
 \mu_s(\Syr_{\min}>K)\leq C_1(\log K)^{-c}
\]
for all sufficiently large $K$, uniformly over nonempty blocks.  Bounded
$K\geq2$ is included by increasing $C_1$.
For $s\geq2$, integral comparison gives $h_s\leq C_\alpha\log s$.
If $1<s<2$, the block contains no odd integer, since $s^\alpha<2^\alpha<3$.
Hence every relevant block obeys
\begin{equation}\label{eq:block-harmonic}
 \sum_{\substack{m\in R_s\\\Syr_{\min}(m)>K}}\frac1m
 \leq C_2(\log s)(\log K)^{-c},
\end{equation}
with empty blocks contributing zero.

If $X\leq K$, the sum in \eqref{eq:odd-fixed} is zero, since the orbit
contains its starting value.  Otherwise let $z_i=X^{\alpha^{-i}}$, and
let $J\geq1$ be minimal with $z_J\leq K$.  The blocks
$[z_i,z_i^\alpha]$, $1\leq i\leq J$, cover $[z_J,X]$.
Every bad starting value exceeds $K$, so the remaining prefix contributes
nothing.  Possible duplicated endpoints only increase an upper bound.
Since
\[
 \sum_{i=1}^J\log z_i
 =\log X\sum_{i=1}^J\alpha^{-i}
 \leq\frac{\log X}{\alpha-1},
\]
summing \eqref{eq:block-harmonic} proves \eqref{eq:odd-fixed}.

For the full-domain map, Propositions~\ref{prop:clocks} and
\ref{prop:dyadic} give the finite identity
\[
 \sum_{\substack{n\leq X\\\Col_{\min}(n)>K}}\frac1n
 =\sum_{a\geq0}2^{-a}
   \sum_{\substack{m\leq X/2^a,\ m\in\Opos\\\Syr_{\min}(m)>K}}\frac1m.
\]
Only finitely many terms are nonzero.  If $X/2^a<2$, the inner bad sum is
zero.  For the others, apply \eqref{eq:odd-fixed} and use
\[
 \sum_{a:X/2^a\geq2}2^{-a}\log(X/2^a)
 \leq\log X\sum_{a\geq0}2^{-a}=2\log X.
\]
This proves \eqref{eq:full-fixed}.
\end{proof}

\subsection{Diverging thresholds and the exact meaning of tightness}

\begin{corollary}[Tao's almost-bounded theorem]\label{cor:tao}
For every $f:\N_+\to\R$ with $f(n)\to+\infty$,
\[
 \lim_{X\to\infty}\frac1{H(X)}
 \sum_{\substack{n\leq X\\\Col_{\min}(n)\geq f(n)}}\frac1n=0.
\]
Thus $\Col_{\min}(n)<f(n)$ on a set of logarithmic density one.
The analogous statement on $\Opos$, with $H_o(X)$ and $\Syr$, also holds.
\end{corollary}
\begin{proof}
Fix an integer $K\geq2$.  There is $Y_K$ such that $f(n)>K$ for
$n\geq Y_K$.  Therefore
\[
 \{n:\Col_{\min}(n)\geq f(n)\}
 \subseteq[1,Y_K]\cup\{n:\Col_{\min}(n)>K\}.
\]
Theorem~\ref{thm:fixed} bounds its harmonic mass up to $X$ by
$H(Y_K)+C\log X/(\log K)^c$.  Divide by $H(X)$ and first let
$X\to\infty$ at fixed $K$.  The resulting limsup is at most
$C/(\log K)^c$, which tends to zero as $K\to\infty$.
The odd case is identical, using $H_o(X)\sim\frac12\log X$.
The argument allows nonmonotone $f$ with arbitrarily slow divergence.
Its negative values occur only in a finite prefix, and the equality part
of the bad event is included throughout.
\end{proof}

The finite-prefix argument gives a general tightness criterion for
weighted counting measures.

\begin{proposition}[Weighted tightness equivalence]\label{prop:tightness}
Let $D\subseteq\N_+$ be infinite, let $w(n)>0$, and suppose
$W(X)=\sum_{n\in D,n\leq X}w(n)\to\infty$.
Let $g:D\to[0,\infty)$ be finite-valued and let $\rho_X$ be the
probability proportional to $w$ on the nonempty prefix $D\cap[1,X]$.
The following assertions are equivalent:
\begin{enumerate}
\item For every real-valued $f(n)\to\infty$, $\rho_X(g\geq f)\to0$.
\item $\displaystyle\lim_{K\to\infty}\limsup_{X\to\infty}\rho_X(g>K)=0$.
\item $\displaystyle\lim_{K\to\infty}\sup_X\rho_X(g>K)=0$.
\end{enumerate}
The supremum is over nonempty prefixes.
\end{proposition}
\begin{proof}
The finite-prefix argument in Corollary~\ref{cor:tao}, with $W$ in place
of $H$, proves (2)$\Rightarrow$(1).
If (2) fails, choose $\varepsilon>0$ such that for every integer $j\geq1$
there are arbitrarily large $X$ with $\rho_X(g>j)\geq\varepsilon$.
Start with an integer $X_0\geq\min D$.
Choose integers $X_j>X_{j-1}$ recursively with this property and with
$W(X_{j-1})/W(X_j)\leq\varepsilon/2$.  Set $f(n)=j$ for
$X_{j-1}<n\leq X_j$, assigning any finite value on the initial prefix.
This tends to infinity.  At $X_j$, removal of the earlier prefix leaves
mass at least $\varepsilon/2$ on which $g>j=f$, contradicting (1).

To prove (2)$\Rightarrow$(3), choose $K_0$ so the limsup in (2) is less
than $\varepsilon/2$.  For some $X_0$, all $X\geq X_0$ have
$\rho_X(g>K_0)\leq\varepsilon$.  Increase $K_0$ to exceed every value
of $g$ on the finite prefix up to $X_0$.  All earlier tails are now empty,
and the later bound persists.  Finally (3)$\Rightarrow$(2) is immediate.
\end{proof}

For $w(n)=1/n$ and $g=\Col_{\min}$ or $\Syr_{\min}$, this gives the
qualitative equivalence stated by Tao.  The quantitative rate in
Theorem~\ref{thm:fixed} uses his first-passage estimate in addition to
this equivalence.
